%==============================================================================
% Research Methods - Groepsopdracht - onderzoeksvoorstellen
%==============================================================================

\documentclass{hogent-article}

% Bibliografiebestand
\addbibresource{references.bib}

% Info over de opdracht, het vak en de opleiding
\studyprogramme{Professionele bachelor toegepaste informatica}
\course{Research Methods, 2E1, 2D}
\assignmenttype{Groepsopdracht}
\academicyear{2025-2026}
\title{Schriftelijke voorbereiding onderzoeksvoorstellen}

\author{Simon Boey}
\email{simon.boey@student.hogent.be}
\sloppy{}
\author{Simon Borghart}
\email{simon.borghart@student.hogent.be}

\author{Emiel De Pelsmaeker}
\email{emiel.depelsmaeker@student.hogent.be}

\author{Milan Saric}
\email{milan.saric@student.hogent.be}

\projectrepo{https://github.com/SaZoMi/RM-Group-64-interprofessionele}

\specialisation{AI \& Data Engineering}

\keywords{deepfake spraakdetectie, AI-stemfraude, vishing, telecombeveiliging,
  lichtgewicht machinaal leren, MFCC, CNN, AVG, privacybehoudende AI}

\begin{document}

\begin{abstract}
  Wanneer een telefonische bankmedewerker om uw rekeninggegevens vraagt, vertrouwt
  u op de menselijkheid van de stem aan de andere kant. Precies dat vertrouwen
  benutten aanvallers die gebruikmaken van synthetische spraak gegenereerd door
  generatieve AI\@. Dit onderzoek evalueert de haalbaarheid van lichtgewicht
  machinaal-leermodellen voor de detectie van AI-gegenereerde frauduleuze
  spraakoproepen onder realistische telecommunicatieomstandigheden. Concreet
  worden CNN-gebaseerde classificatoren getraind op MFCC- en
  mel-spectrogramrepresentaties toegepast op korte audiosegmenten van 2 tot 5
  seconden. Drie onderzoeksvragen staan centraal: in welke mate kunnen
  lichtgewicht modellen synthetische spraak detecteren vanuit korte fragmenten,
  hoe degradeert de detectieprestatie onder codeccompressie en achtergrondlawaai,
  en welke afruil bestaat er tussen detectienauwkeurigheid en
  AVG-conforme privacybeperkingen zoals on-device inferentie. Met behulp van
  de publiek beschikbare ASVspoof~2019-dataset, aangevuld met gesimuleerde
  telecomdegradatie, wordt een proof-of-concept-pijplijn ontwikkeld die
  basismodellen vergelijkt met een lichtgewicht CNN\@. De verwachte uitkomst is
  een gekwantificeerd inzicht in de detectiehaalba\-arheid en haar grenzen onder
  realistische telecom- en privacybeperkingen, wat een empirische basis legt voor
  toekomstige privacybehoudende fraudedetectiesystemen in de Europese
  telecomsector.
\end{abstract}

\tableofcontents

%------------------------------------------------------------------------------
\section{Inleiding}%
\label{sec:inleiding}
%------------------------------------------------------------------------------

Wie vertrouwt u meer: een onbekend telefoonnummer op uw scherm, of de herkenbare
stem van uw bankdirecteur? Stemherkenning is voor mensen een diepgeworteld
vertrouwenssignaal, en dat maakt het bijzonder kwetsbaar voor misbruik.
Hedendaagse spraaksynthesesystemen zijn inmiddels in staat om een overtuigende
kloon van een doelspreekstem te produceren op basis van slechts enkele seconden
referentie-audio~\autocite{Azzuni2025}. De drempel voor het uitvoeren van
zogeheten \emph{vishing}-aanvallen (voice phishing) is derhalve dramatisch
gedaald: wat vroeger gespecialiseerde apparatuur vereiste, is vandaag uitvoerbaar
via openbaar beschikbare cloudAPIs en open-source bibliotheken~\autocite{Ruggiero2021}.

Bestaande fraudedetectiesystemen bij telecomoperatoren steunen doorgaans op
metadataanalyse: oproepfrequentie, nummerreputatie en gedragsafwijkingen op
abonnementsniveau. Deze aanpak volstaat voor grootschalige spamcampagnes, maar
schiet structureel tekort bij gerichte AI-gegenereerde spraakaanvallen waarbij
elke oproep zich op netwerkniveau als volledig legitiem voordoet. De aanvalsvector
verschuift bijgevolg van de metadatalaag naar de inhoud van de audio zelf.

Anderzijds heeft onderzoek naar de detectie van synthetische spraak via machinaal
leren aanzienlijke vooruitgang geboekt, maar doorgaans onder veronderstellingen
die in productieomgevingen niet opgaan: volledige opnames van hoge kwaliteit,
geen latentiebeperkingen en gecentraliseerde cloudverwerking. In de Europese
context (en specifiek in België) vormt de Algemene Verordening
Gegevensbescherming (AVG) een extra regulatoire barrière~\autocite{GDPR2016}:
stemdata wordt geclassificeerd als persoonsgegevens en potentieel als biometrische
gegevens, wat de opslag en verzending ervan naar centrale servers sterk
beperkt~\autocite{Nautsch2019}.

Hieruit volgt een concrete lacune: succesvolle laboratoriummodellen voor
deepfake-spraakdetectie zijn moeilijk of niet inzetbaar in
real-time, privacybehoudende telecominfrastructuur. Dit onderzoek beoogt die
lacune empirisch in kaart te brengen.

\subsection*{Onderzoeksdoelstelling}

Het doel van dit onderzoek is de haalbaarheid te evalueren van lichtgewicht
machinaal-leermodellen voor de detectie van AI-gegenereerde of fraudegerelateerde
spraakpatronen onder realistische telecoombeperkingen, met behulp van een
proof-of-concept-aanpak.

\subsection*{Onderzoeksvragen}

\begin{description}
  \item[OV1] In welke mate kunnen lichtgewicht machinaal-leermodellen
    AI-gegenereerde of frauduleuze spraak detecteren op basis van uitsluitend
    korte audiosegmenten van 2 tot 5 seconden?
  \item[OV2] Hoe degradeert de detectieprestatie onder realistische
    telecommunicatieomstandigheden zoals codeccompressie, achtergrondlawaai en
    onvolledige spraakfragmenten?
  \item[OV3] Welke afruil bestaat er tussen detectieprestatie en
    privacybehoudende beperkingen zoals on-device verwerking en het vermijden van
    audio-opslag of -transmissie?
\end{description}

%------------------------------------------------------------------------------
\section{Literatuurstudie}%
\label{sec:literatuurstudie}
%------------------------------------------------------------------------------

\subsection{Synthetische spraakgeneratie en stemklonen}

De afgelopen jaren heeft de opmars van neurale spraaksynthese geleid tot
text-to-speech (TTS)- en voice-conversion (VC)-systemen die menselijke prosodie,
accent en emotionele variatie nauwkeurig kunnen nabootsen op basis van beperkte
referentie-audio~\autocite{Azzuni2025}. Transferlearning-benaderingen maken het
voorts mogelijk om TTS-modellen te generaliseren naar nieuwe sprekers zonder
sprekerspecifieke trainingsdata, waardoor de productie van overtuigende
stemklonen geen gespecialiseerde expertise meer vereist~\autocite{Ruggiero2021}.
De brede beschikbaarheid van dergelijke systemen via opensourceplatformen heeft
de drempel voor misbruik in fraudescenario's, waaronder identiteitsdiefstal,
financiële oplichting en desinformatiecampagnes, drastisch verlaagd.

\subsection{Detectiemethoden voor audio-deepfakes}

Onderzoek naar de detectie van synthetische spraak wordt grotendeels gestructureerd
rond de ASVspoof-challengereeks, die publiek beschikbare benchmarkdatasets
aanbiedt van authentieke en gespoofde spraakfragmenten~\autocite{Wang2020}. De
ASVspoof~2019 Logical Access-partitie is uitgegroeid tot het referentiecorpus voor
TTS- en VC-detectie, met fragmenten afkomstig van 19 verschillende
synthesemethoden die een breed spectrum aan kwaliteit en sprekersgelijkenis
vertegenwoordigen.

In de literatuur hebben zich twee hoofdfamilies van detectorarchitecturen
afgetekend. Enerzijds steunen op maat gemaakte feature-gebaseerde methoden op
representaties zoals Mel-Frequency Cepstral Coefficients (MFCC's),
Linear-Frequency Cepstral Coefficients (LFCC's) of spectrale statistieken,
gekoppeld aan klassieke classificatoren. Anderzijds passen deep-learningmodellen
convolutionele neurale netwerken (CNN's) of transformerarchitecturen rechtstreeks
toe op tijd-frequentierepresentaties zoals mel-spectrogrammen~\autocite{Bartusiak2022}.
Ensemblemodellen en zelfgesuperviseerde feature-extractoren hebben op de
ASVspoof-benchmarks sterke resultaten laten optekenen, met Equal Error Rates
(EER) van minder dan 1\% onder gecontroleerde omstandigheden.

\subsection{Beperkingen van bestaande detectiesystemen}

Ondanks indrukwekkende benchmarkresultaten stoten bestaande deepfake-detectoren
op fundamentele beperkingen zodra ze buiten gecontroleerde laboratoriumomgevingen
worden ingezet. Een centrale probleemfactor is de gevoeligheid voor
audiodegradatie: codeccompressie, achtergrondlawaai, bandbreedtebeperking,
pakketverlies en resampling ondermijnen de modelprestaties
merkbaar~\autocite{Shi2025}. \textcite{Li2025} evalueerden tien
state-of-the-art-detectoren systematisch aan de hand van 16 reële
audiocorrupties en stelden vast dat, hoewel de meeste modellen additief lawaai
redelijk goed doorstaan, ze sterk kwetsbaar blijven voor compressie-artefacten, met name van neurale
codecs die de laagniveau-spectrale kenmerken vernietigen waarop detectoren steunen.

In aansluiting hierop toonden \textcite{Yadav2024} aan dat audio van
telefoongesprekskwaliteit de prestaties van standaarddetectoren aanzienlijk
degradeert, en stelden een transformerarchitectuur voor die specifiek is ontworpen
voor gecomprimeerde invoer. Nochtans blijft generalisatie een
hardnekkig probleem: modellen die op één set synthesemethoden zijn getraind,
presteren doorgaans zwak op nieuwe TTS- of VC-systemen die tijdens de training
niet aanwezig waren.

\subsection{Telecomfraude en de detectie van oplichtersgesprekken}

Fraudedetectie in telecommunicatienetwerken heeft van oudsher gesteund op
netwerklevel-metadata: oproepfrequentie, nummerreputatie en
abonnementsgedragsafwijkingen. Dergelijke benaderingen zijn effectief tegen
massale robocalling-campagnes, maar schieten tekort bij gerichte AI-stem\-fraude,
waarbij elk gesprek op netwerkniveau als legitiem wordt beoordeeld. Deze
verschuiving in het dreigingsmodel impliceert dat de audio-inhoud zelf als
primair fraudesignaal moet fungeren in plaats van uitsluitend oproepmetadata.
De integratie van audio-inhoudsanalyse in telecominfrastructuur introduceert
evenwel aanzienlijke technische uitdagingen met betrekking tot real-time
latentie, schaalbaarheid op operatorniveau en regulatoire naleving.

\subsection{Privacy en AVG-beperkingen bij spraakverwerking}

Op grond van de AVG~\autocite{GDPR2016} vormen spraakopnames persoonsgegevens
en kunnen zij als biometrische gegevens worden geclassificeerd wanneer ze worden
verwerkt met het oog op de unieke identificatie van een natuurlijk persoon.
\textcite{Nautsch2019} wijzen op een substantiële kloof tussen het juridische
kader en de gangbare praktijk in spraaktechnologie: de meeste ingezette systemen
zenden ruwe audio door naar gecentraliseerde cloudinfrastructuur, wat in
tegenspraak kan zijn met de AVG-beginselen van dataminimalisering,
doelbinding en opslagbeperking.

Artikel~9 van de AVG legt verstrekkende verwerkingsbeperkingen op voor
biometrische gegevens, waaronder de vereiste van uitdrukkelijke toestemming
of een specifieke wettelijke grondslag. In de context van telecom\-fraude\-detectie
dwingt dit tot privacybehoudende inferentie-architecturen waarbij ruwe audio
noch wordt opgeslagen noch wordt verzonden, een vereiste die on-device of
edge-gebaseerde verwerking als voornaamste technische strategie aanwijst voor
AVG-conforme inzetting in Europa.

\subsection{Lichtgewicht machinaal leren voor audioclassificatie}

Vanwege de rekenintensiviteit van grootschalige deep-learningmodellen is er een
groeiende onderzoeksstroom gericht op lichtgewicht architecturen die geschikt zijn
voor real-time inferentie op randapparaten. \textcite{Amiriparian2022} toonden
aan dat CNN's met beeldpre-training, toegepast op on-the-fly gegenereerde
mel-spectrogrammen, real-time inferentie kunnen uitvoeren op consumentsmartphones
met een gemiddelde latentie van minder dan 250~ms. \textcite{Bartusiak2022}
toonden voorts aan dat compacte convolutionele transformers die op
spectrogrampatches werken, competitief zijn met beduidend grotere architecturen
voor synthetische spraakdetectie.

De fundamentele afruil in dit domein betreft modelomvang versus
generalisatiecapaciteit: kleinere modellen kunnen onderpresteren op gedegradeerde
of out-of-distribution audio die kenmerkend is voor reëel telecomverkeer.
Het kwantificeren van die afruil onder realistische omstandigheden vormt een
centrale motivatie voor dit onderzoek.

\subsection{Samenvatting van de onderzoekslacune}

Samenvattend heeft de literatuur sterke vooruitgang geboekt in deepfake-audio\-detectie
onder gecontroleerde omstandigheden, maar weinig onderzoek richt zich specifiek
op de gelijktijdige aanwezigheid van drie beperkingen die kenmerkend zijn voor
reële telecomsystemen: (1) korte detectievensters van slechts 2 tot 5~seconden
voor vroegtijdige analyse van oproepen, (2) realistische audiodegradatie door
codeccompressie en pakketverlies, en (3) AVG-conforme on-device inferentie zonder
persistente audio-opslag. De combinatie van deze drie beperkingen is rechtstreeks
relevant voor operationele inzetting in Europese telecomnetwerken en definieert
de lacune die dit onderzoek beoogt te adresseren.

%------------------------------------------------------------------------------
\section{Methodologie}%
\label{sec:methodologie}
%------------------------------------------------------------------------------

Dit onderzoek volgt een proof-of-concept experimentele aanpak, gestructureerd
in vier fasen: datavoorbereiding, feature-extractie, modellering en evaluatie.

\subsection{Datavoorbereiding}

De primaire dataset is de ASVspoof~2019 Logical Access-partitie~\autocite{Wang2020},
die authentieke spraakfragmenten en synthetische samples biedt van 19
verschillende TTS- en VC-systemen. Aanvullende synthetische samples worden
gegenereerd via publiek beschikbare TTS-systemen om de variëteit aan aanvalstypes
te verbreden. Alle audio wordt gesegmenteerd in vaste vensters van 2, 3 en 5~seconden
teneinde de real-world beperking op vroegtijdige oproepanalyse te simuleren.

Om realistische telecomdegradatie te simuleren, worden de volgende augmentaties
toegepast:

\begin{itemize}
  \item codeccompressie via MP3 en Opus op meerdere bitrates (8, 16 en 32~kbps);
  \item injectie van achtergrondlawaai bij signaal-ruisverhouding van 5, 10
    en 20~dB;
  \item gesimuleerd pakketverlies van 5\% en 10\%;
  \item telefoonband-frequentiefiltering (300--3400~Hz) conform GSM-specificaties.
\end{itemize}

\subsection{Feature-extractie}

Twee feature-representaties worden onafhankelijk geëvalueerd:

\begin{itemize}
  \item \textbf{MFCC}: 40 Mel-Frequency Cepstral Coefficients berekend per
    frame met een venster van 25~ms en een hopgrootte van 10~ms, samengevoegd
    tot een vaste matrix per audiosegment via gemiddeld en standaarddeviatiepooling.
  \item \textbf{Mel-spectrogram}: logaritmisch geschaalde mel-filterbankoutput
    met 128 mel-banden, voorgesteld als 2D-afbeelding geschikt voor directe
    CNN-invoer.
\end{itemize}

\subsection{Modellen}

Drie modelfamilies worden getraind en onderling vergeleken:

\begin{enumerate}
  \item \textbf{Basismodellen}: logistische regressie en random forest op
    afgeplatte MFCC-features, als interpreteerbare prestatie-referentie zonder
    deep learning.
  \item \textbf{Lichtgewicht CNN}: een compact convolutioneel netwerk toegepast
    op mel-spectrogramafbeeldingen, ontworpen voor CPU-gebaseerde real-time
    inferentie conform de aanpak van~\textcite{Amiriparian2022}.
  \item \textbf{Compact convolutioneel transformer}: een kleinschalig
    attentiemodel toegepast op spectrogrampatches naar het voorbeeld
    van~\textcite{Bartusiak2022}, als nauwkeurigheids-bovengrens ter vergelijking
    met het lichtgewicht CNN\@.
\end{enumerate}

Alle modellen worden getraind op schone audio en vervolgens geëvalueerd op zowel
schone als gedegradeerde testsets, zodat de generalisatiecapaciteit onder
telecomcondities meetbaar wordt.

\subsection{Evaluatiemetrieken}

De modellen worden beoordeeld op de volgende metrieken:

\begin{itemize}
  \item Precisie, Recall en macro-F1-score;
  \item ROC-AUC;
  \item Equal Error Rate (EER), conform de ASVspoof-evaluatieconventie;
  \item \emph{Robuustheidsdelta}: absolute F1-degradatie tussen schone en
    maximaal gedegradeerde testcondities;
  \item inferentielatentie in milliseconden per 5-secondenvenster, gemeten op
    een CPU-only host ter simulatie van edge-inzetting.
\end{itemize}

%------------------------------------------------------------------------------
\section{Verwachte resultaten}%
\label{sec:verwachte-resultaten}
%------------------------------------------------------------------------------

Dit onderzoek beoogt een werkende proof-of-concept-pijplijn op te leveren die de
haalbaarheid en de grenzen van lichtgewicht machinaal leren voor deepfake-spraak\-detectie
onder telecoombeperkingen aantoont. Concreet worden de volgende uitkomsten
verwacht:

\begin{itemize}
  \item Een functionele binaire classificatiepijplijn die synthetische van
    authentieke spraak onderscheidt met aanvaardbare nauwkeurigheid op schone
    audio, als vertrekpunt voor verder onderzoek.
  \item Gekwantificeerde prestatiedegradatiecurves als functie van codecbitrate,
    ruisniveau en vensterlengte, die de operationele drempelwaarden identificeren
    waaronder detectie onbetrouwbaar wordt.
  \item Een vergelijkende analyse van modelomvang versus detectieprestatie en
    inferentielatentie, die vaststelt of lichtgewicht CNN's een haalbare
    prestatie-efficiëntie-afruil bieden voor on-device inzetting.
  \item Empirische inzichten in de minimale audioduur die betrouwbare
    vroegtijdige detectie mogelijk maakt onder realistische degradatie.
\end{itemize}

Op grond van de bevindingen van~\textcite{Li2025} en~\textcite{Shi2025}
wordt verwacht dat de detectieprestatie aanvaardbaar blijft bij schone of licht
gecomprimeerde audio, maar significant degradeert onder zware codeccompressie.
Dit wijst op de noodzaak van compressiebewuste trainingstrategieën of
data-augmentatie voor elke praktische inzetting. Het lichtgewicht CNN wordt
verondersteld een inferentielatentie ruim onder de 250~ms per segment te
bereiken~\autocite{Amiriparian2022}, wat de real-time haalbaarheid op
edge-hardware bevestigt.

%------------------------------------------------------------------------------
\section{Conclusie}%
\label{sec:conclusie}
%------------------------------------------------------------------------------

Dit onderzoek adresseert een praktisch relevante en empirisch onderbelichte
probleemstelling: de real-time detectie van AI-gegenereerde stemfraude onder de
gecombineerde beperkingen van reële telecommunicatieomgevingen en Europees
gegevensbeschermingsrecht. In tegenstelling tot onderzoek dat streeft naar
maximale nauwkeurigheid in gecontroleerde omgevingen, is de primaire bijdrage
van deze studie een kwantitatieve evaluatie van de mate waarin drie gelijktijdige
real-world beperkingen (korte audiovensters, telecom-codeccompressie en
AVG-conforme privacybeperkingen) de haalbaarheid van lichtgewicht deepfake\-detectoren
ondermijnen.

De resultaten beogen te verduidelijken of bestaande lichtgewicht architecturen
rechtstreeks inzetbaar zijn voor vroegtijdige fraudedetectie, dan wel of gerichte
aanpassingen (zoals compressiebewuste training, on-the-fly data-augmentatie of
gespecialiseerde feature-representaties) noodzakelijke precondities zijn voor
operationele inzetting. Hiermee positioneert de studie zich als een
haalbaarheidsevaluatie die de ontwerpeisen voor elk productiereed systeem
empirisch onderbouwt.

Praktische beperkingen van het onderzoek omvatten de afhankelijkheid van
gesimuleerde in plaats van authentieke oproepdata en de gecontroleerde aard van
de telecomdegradatiesimulatie, die de volledige diversiteit van reële netwerk\-omstandigheden
mogelijk niet omvat. Toekomstige onderzoeksrichtingen omvatten multimodale
fraudedetectie waarbij audio-analyse wordt gecombineerd met oproepmetadata,
federatief of volledig on-device leren voor AVG-conforme modeltraining, integratie
in reële telecomoperatorinfrastructuur onder gecontroleerde pilootomstandigheden,
en de evaluatie van meer geavanceerde deepfake-architecturen zoals die voortdurend
worden ontwikkeld in de ASVspoof-onderzoeksgemeenschap.

\printbibliography[heading=bibintoc]

\end{document}
